%=============================================================================
%  03_product.tex  -  Chapter 2: Product Philosophy and Users
%=============================================================================
\chapter{Product Philosophy and Users}
\label{ch:product}

\section{Who it is for}
Universal Analyst is built for \textbf{data analysts, data-science students, and
technical operators} who arrive with a concrete dataset and a concrete question.
Their context is a focused analysis session rather than a long-lived project:
they upload data, read the profile, ask grounded questions, request a chart or a
prediction, and export a report. They value \emph{speed}, \emph{density}, and
\emph{trustworthy output} over decoration, and they are fluent in tools like
Linear and Stripe --- so standard product affordances should behave exactly as
they expect.

\begin{table}[h]
\centering
\caption{Primary persona and the job the product does for them.}
\begin{tabularx}{\linewidth}{@{}L{3.4cm} X@{}}
\toprule
\textbf{Persona} & \textbf{Job-to-be-done} \\
\midrule
Data analyst & ``I have this file and a question --- give me a correct,
grounded answer and, if it helps, a real chart, without me writing code.'' \\
DS student & ``Let me explore a dataset's structure, try a quick predictive
model, and produce a clean report I can hand in.'' \\
Technical operator & ``Run this locally on our GPU box; nothing about our data
should leave the machine.'' \\
\bottomrule
\end{tabularx}
\end{table}

\section{Brand personality}
The product is deliberately positioned as an \emph{instrument}, not a marketing
app: precise, technical, and quietly confident. Three words anchor the identity
--- \textbf{measured, grounded, sharp}. It should feel like a piece of lab
equipment a professional reaches for, where every readout is exact and earned:
no hype, no hand-holding, no filler copy.

\section{Design principles}
The following principles govern both the interface and the backend behaviour.
They are not aspirational --- each maps to concrete implementation decisions
described later in this document.

\begin{description}[style=nextline,leftmargin=0pt]
  \item[\textcolor{uaIndigoD}{Readouts are exact.}]
    Numbers, counts, and labels render in a monospaced face (IBM Plex Mono) so
    data reads like instrument output rather than prose. See the front-end
    typography in Chapter~\ref{ch:frontend}.
  \item[\textcolor{uaIndigoD}{Earned familiarity.}]
    Tabs, side navigation, the composer, and tables behave exactly as a
    Linear/Stripe-fluent user expects, so the tool disappears into the task.
  \item[\textcolor{uaIndigoD}{No fabrication.}]
    The UI never substitutes static or placeholder analysis for real model
    output. Empty and error states say what actually happened --- this is
    enforced in the streaming chat path (Chapter~\ref{ch:llm}).
  \item[\textcolor{uaIndigoD}{One signal, one job.}]
    Indigo carries actions, focus, and identity; amber marks measurement and
    emphasis only. Neither colour is used to decorate.
  \item[\textcolor{uaIndigoD}{Density with rhythm.}]
    The interface is information-dense where the analyst needs it, using
    deliberate spacing and hairline structure rather than nested cards.
\end{description}

\begin{uanote}[The no-fabrication rule, concretely]
When the model or Ollama fails mid-answer, the backend emits a real
\code{error} event and the front end shows it verbatim. The system never
back-fills a canned profile summary, a fake chart, or invented statistics to
paper over a failure. This single rule shapes error handling across every
subsystem.
\end{uanote}

\section{Anti-references}
Part of defining a product is defining what it must \emph{not} become. Universal
Analyst explicitly rejects:

\begin{itemize}
  \item The generic AI-SaaS dashboard --- rounded cards everywhere, purple-to-blue
        gradient accents, a tiny uppercase eyebrow over every section, and
        hero-metric templates.
  \item The ``warm cream and terracotta serif'' editorial template (the
        project's original look, explicitly retired).
  \item The flat ``near-black plus one acid-green accent'' dark default; the
        current theme stays differentiated by its indigo/amber duotone and
        monospaced data signature.
  \item Any surface that advertises a capability it does not actually have.
\end{itemize}

\section{Accessibility and inclusion}
The application ships a dark theme in which body text and labels meet WCAG AA
contrast (at least \code{4.5:1}, or \code{3:1} for large text). Every
interactive element has a visible keyboard focus state, all motion has a
\code{prefers-reduced-motion} fallback, and microphone dictation is an optional
enhancement --- never the only path to input. Accessibility is treated as a
correctness property of the interface, not a later polish step.
